\notechapter{N-20221228}{Quotient Topology and Product Topology}

\section{Interval modulo endpoints}

Consider the quotient
\[
  D^1=[-1,1],\qquad -1\sim1.
\]
Identifying the two endpoints gives a circle:
\[
  D^1/(-1\sim1)\cong S^1.
\]
A small neighborhood around the glued point corresponds to two half-interval neighborhoods at \(-1\) and \(1\) upstairs.  This is exactly the local reason quotient topology is defined by preimages.

More generally,
\[
  D^n/S^{n-1}\cong S^n.
\]
Collapsing the boundary of a disk to one point produces a sphere.

\section{More quotient spaces}

The standard square identifications are determined by the direction of edge gluing.  Identifying opposite vertical sides in the same direction gives a cylinder; identifying the two circular boundary components of that cylinder in the same direction gives a torus.  A twist in one pair of edge identifications produces a Möbius band or a Klein bottle depending on the full pattern.

Another example is
\[
  \mathbb R/\mathbb Z\cong S^1,
\]
where \(x\sim y\) iff \(y-x\in\mathbb Z\).  This is the real line wrapped around the circle.

\section{Definition of quotient topology}

\begin{definition}[Quotient topology]
Given a surjective map \(q:X\to Y\), the quotient topology on \(Y\) declares \(U\subseteq Y\) open exactly when \(q^{-1}(U)\) is open in \(X\).
\end{definition}


Let \(q:X\to X/{\sim}\) be the quotient map.  A set \(U\subseteq X/{\sim}\) is open iff
\[
  q^{-1}(U)\text{ is open in }X.
\]
This is the finest topology on the quotient making \(q\) continuous.  The note emphasizes that quotient topology is not guessed by drawing; it is forced by the universal property of the quotient map.

\section{Product topology}

For spaces \(X,Y\), the product topology on \(X\times Y\) is generated by basic open rectangles
\[
  U\times V,\qquad U\subseteq X\text{ open},\quad V\subseteq Y\text{ open}.
\]
Thus a set is open in \(X\times Y\) if it is a union of such rectangles.

The page's picture of a small rectangle in a grid illustrates the basic rule: local neighborhoods in a product are products of local neighborhoods in the factors.


\section{Universal property of quotient topology}

The quotient topology is not chosen by visual intuition alone.  Let \(q:X\to Y=X/{\sim}\) be the quotient map.  A map \(f:Y\to Z\) is continuous iff
\[
  f\circ q:X\to Z
\]
is continuous.  This is the universal property of the quotient topology.

This property explains the definition.  We declare \(U\subseteq Y\) open precisely when \(q^{-1}(U)\) is open in \(X\), because this is exactly what makes continuity out of the quotient detectable before quotienting.

\section{Endpoint identification in detail}

For \([0,1]/(0\sim1)\), points near the glued point come from two pieces upstairs: a small interval near \(0\) and a small interval near \(1\).  This is why the quotient looks like a circle rather than an interval with a strange endpoint.  The local neighborhoods have been glued.

Similarly, \(D^2/S^1\cong S^2\): collapse the boundary circle of a disk to a single point.  The interior of the disk becomes the sphere with one point removed, and the collapsed boundary becomes the missing point.

\section{Product topology via projections}

The product topology on \(X\times Y\) is the coarsest topology making the projection maps
\[
  \pi_X:X\times Y\to X,
  \qquad
  \pi_Y:X\times Y\to Y
\]
continuous.  Its basis consists of rectangles
\[
  U\times V
\]
where \(U\subseteq X\) and \(V\subseteq Y\) are open.

The word “coarsest” matters.  We do not add more open sets than necessary.  Once all projection preimages \(\pi_X^{-1}(U)=U\times Y\) and \(\pi_Y^{-1}(V)=X\times V\) are open, finite intersections give \(U\times V\), and arbitrary unions give the full product topology.

\section{Quotient and product as opposite directions}

Product topology builds a larger space from two spaces while preserving maps out to the factors.  Quotient topology builds a smaller space by identifying points while preserving maps out of the quotient.  Product is controlled by projections; quotient is controlled by the quotient map.

This contrast is a useful memory:
\[
  \text{product: maps out to factors},\qquad
  \text{quotient: maps out after gluing}.
\]
Both definitions are forced by continuity requirements rather than by pictures.

\section{A second lens}

Quotient topology and product topology are dual-looking constructions.  A quotient identifies points and defines open sets by pulling back along the quotient map.  A product combines spaces and defines open sets by rectangles from the two factors.  Both are examples of topology being controlled by maps rather than by drawings alone.

\section{Saturated sets and quotient openness}

For a quotient map \(q:X\to Y\), a subset \(A\subseteq X\) is saturated if it is a union of equivalence classes, equivalently
\[
  A=q^{-1}(q(A)).
\]
Open sets in \(Y\) correspond exactly to saturated open sets in \(X\).  This is a practical way to test quotient topology: first check that the upstairs set is open, then check that it contains whole equivalence classes.

Endpoint gluing works because neighborhoods near the glued point have saturated preimages containing pieces near both endpoints.  A one-sided interval near only one endpoint is not saturated, so it does not define a neighborhood downstairs.

\section{Quotients can destroy Hausdorffness}

Not every quotient of a nice space is nice.  For example, identify all rational numbers in \(\mathbb R\) to one point and all irrational numbers to another point.  The quotient has two points, but the preimage of either singleton is dense in \(\mathbb R\), so the two quotient points cannot be separated by disjoint open neighborhoods.

This warning matters because quotient pictures can be misleading.  The quotient topology is defined by preimages, and sometimes that produces spaces with poor separation properties.

\section{Product topology and the mapping property}

The product topology is characterized by a universal property dual to the quotient one.  A map
\[
  h:Z\to X\times Y
\]
is continuous iff both coordinate maps
\[
  \pi_X\circ h:Z\to X,\qquad \pi_Y\circ h:Z\to Y
\]
are continuous.  This is why the product topology is the coarsest topology making the projections continuous.

Thus product and quotient are both controlled by maps:
\[
\begin{array}{c|c|c}
\text{construction} & \text{distinguished map} & \text{continuity test}\\
\hline
\text{quotient} & q:X\to X/{\sim} & f\text{ out of quotient via }f\circ q\\
\text{product} & \pi_X,\pi_Y & h\text{ into product via coordinates}
\end{array}
\]

\section{From squares to surfaces through edge words}

The quotient square diagrams can be encoded algebraically.  A torus has boundary word
\[
  aba^{-1}b^{-1},
\]
while a Klein bottle has a word of the form
\[
  aba^{-1}b.
\]
The arrows in the picture determine whether a generator appears as \(a\) or \(a^{-1}\).  This is already the language of CW complexes and fundamental-group presentations.

So the square gluing is not merely visual.  It tells us how a two-cell is attached to a one-skeleton, and the attaching word becomes a relation in the fundamental group.

\section{Products and quotients need not commute blindly}

Although \((\mathbb R/\mathbb Z)\times(\mathbb R/\mathbb Z)\) gives the torus, one should not assume products and quotients always interact without hypotheses.  The map
\[
  X\times Y\to (X/{\sim})\times(Y/{\approx})
\]
passes through the quotient by the product equivalence relation, but whether the resulting topology behaves as expected depends on the quotient maps and spaces involved.

This is another reason universal properties are safer than pictures: they tell us which maps are actually continuous and which topology is being imposed.

\section{Products, quotients, and gluing data}

Quotient topology is about gluing while controlling maps out of the glued space.  Product topology is about pairing spaces while controlling maps into the product.  Endpoint identification, \(\mathbb R/\mathbb Z\), torus gluings, and square diagrams are not separate examples; they are the same principle expressed through equivalence classes, saturated sets, and universal properties.
